\documentclass[12pt]{article}
\usepackage[utf8]{inputenc}
\usepackage[T1]{fontenc}
\usepackage{hyperref}
\title{Netvaar Language Resources for K-12 Curriculum Development}
\author{Compiled by Claude Sonnet 5 (claude-sonnet-5) for Andr\'as Kornai}
\date{}
\begin{document}
\maketitle

\section{Introduction}
Netvaar (ISO 639-3 tnl; Glottocode lena1238) is the native name of the language usually called Lenakel or West Tanna. It is an Austronesian (Oceanic) language spoken on the west coast of Tanna Island, Vanuatu. Estimates of the number of speakers range from about 7,000 to 12,000. The search found no school curriculum materials in the language, and no evidence that it is taught in schools. This dossier is therefore short, and the resources below are mostly reference works.

\section{Primary Educational Resources}
No primers, textbooks or school programmes were found in the search results. Searching under the name ``Lenakel'' is more productive than searching under ``Netvaar''.

\section{Linguistic and Reference Materials}
\begin{itemize}
\item Lenakel language (encyclopedia overview): \url{https://en.wikipedia.org/wiki/Lenakel_language}
\item Omniglot, Lenakel (Netvaar), with notes on writing and links: \url{https://www.omniglot.com/writing/lenakel.htm}
\item Languages of Tanna (overview of the island's languages): \url{https://nthieberger.net/vanlangs/TANNA.HTM}
\item The Cambridge Language Centre page on the Tanna languages: \url{https://www.langcen.cam.ac.uk/resources/langt/tanna.html}
\item Lenakel in the Tanna Law, Land and Order (TLOPO) language database: \url{https://tlopo.clld.org/languages/lenakel}
\end{itemize}
John Lynch's \emph{Lenakel Dictionary} (1977) and \emph{A Grammar of Lenakel} (1978) are reported in the search summaries as the main descriptive works. No URL for either was seen in the results, so none is given.

\section{Religious and Community Resources}
\begin{itemize}
\item Scripture Earth page for tnl: \url{https://scriptureearth.org/00eng.php?iso=tnl}
\item Joshua Project language profile: \url{https://joshuaproject.net/index.php/languages/tnl}
\end{itemize}
These sources indicate that a Lenakel New Testament exists (an edition dated 1902 is mentioned), and that there is no complete Bible. I did not verify the text.

\section{Digital and Audio Resources}
\begin{itemize}
\item Linguisten.de note on the LDC Australia and Oceania Lenakel recordings: \url{https://blog.linguisten.de/post/628683434609475585/ldc-australia-oceania-17-lenakel}
\end{itemize}
The Swarthmore Talking Dictionary is reported to have a Lenakel entry (talkingdictionary.swarthmore.edu/lenakel). That address appeared only in text, not as a search-result link, so it has not been checked.

\section{Curriculum Development Context}
Little is available for school use. A curriculum would have to start from the Lynch dictionary and grammar, the New Testament text, and community-led work, none of which was found online here. Local authorities and speakers should be asked what orthography is in use.
\end{document}